\section{Mở rộng lý thuyết: Self-Attention và Transformer}

Mặc dù GRU đạt kết quả khả quan trên IMDB, kiến trúc hồi quy có hai nhược điểm cấu trúc không thể khắc phục bằng việc điều chỉnh siêu tham số. Thứ nhất, phép cập nhật $h_t = (1 - z_t) \odot n_t + z_t \odot h_{t-1}$ buộc phải thực hiện tuần tự trên trục thời gian: muốn tính $h_T$ phải tính trước $h_1, h_2, \dots, h_{T-1}$, dẫn tới độ sâu \textit{wall-clock} của \textit{forward pass} là $\Theta(T)$ ngay cả trên phần cứng song song hóa tốt như GPU. Thứ hai, đường đi của thông tin từ token thứ $i$ tới đầu ra phải đi qua $T-i$ phép cập nhật có cổng; dù cơ chế cổng giảm thiểu hiện tượng \textit{vanishing gradient} so với RNN cổ điển, mỗi bước vẫn ghi đè một phần thông tin, khiến các quan hệ phụ thuộc xa dần bị làm mờ.

Cơ chế \textit{self-attention} (Vaswani và cộng sự, 2017) khắc phục cả hai nhược điểm trên bằng cách thay thế hồi quy bằng một phép \textit{pooling} toàn cục song song. Với $X \in \mathbb{R}^{T \times d}$ là ma trận biểu diễn của chuỗi đầu vào, phép chú ý có dạng:
\begin{equation}
\text{Attention}(Q, K, V) = \text{softmax}\left( \frac{Q K^\top}{\sqrt{d_k}} \right) V
\end{equation}
với $Q = X W_Q$, $K = X W_K$, $V = X W_V$ là các phép chiếu học được vào không gian truy vấn (\textit{query}), khóa (\textit{key}) và giá trị (\textit{value}). Hệ số $1/\sqrt{d_k}$ ngăn \textit{softmax} bão hòa khi $d_k$ lớn, do phương sai của tích vô hướng tăng tuyến tính theo $d_k$. 

Phiên bản \textit{multi-head} ghép nối $H$ phép chú ý độc lập:
\begin{equation}
\text{MultiHead}(X) = \text{Concat}(\text{head}_1, \dots, \text{head}_H) W_O
\end{equation}
cho phép mỗi \textit{head} học một dạng quan hệ ngữ nghĩa riêng biệt: quan hệ chủ-vị, quan hệ phủ định, đồng tham chiếu, hoặc các mối liên kết dài hạn khác.

Sự khác biệt giữa hai kiến trúc được tổng hợp ở Bảng~\ref{tab:so-sanh-gru-attention}. \textit{Self-attention} đổi chi phí tính toán hồi quy $\mathcal{O}(T \cdot d^2)$ lấy chi phí \textit{self-attention} $\mathcal{O}(T^2 \cdot d)$, một sự đánh đổi hoàn toàn chấp nhận được với $T=256$ và $d=128$, đồng thời giải phóng được hai \textit{bottleneck} quan trọng nhất.

\begin{table}[htbp]
\centering
\caption{So sánh GRU và Self-Attention theo các tiêu chí cấu trúc.}
\label{tab:so-sanh-gru-attention}
\begin{tabularx}{\textwidth}{l >{\raggedright\arraybackslash}X >{\raggedright\arraybackslash}X}
\toprule
\textbf{Tiêu chí} & \textbf{GRU} & \textbf{Self-Attention} \\ \midrule
Độ sâu tuần tự (\textit{forward pass}) & $\Theta(T)$ & $\Theta(1)$, song song toàn bộ \\ \addlinespace
Độ dài đường đi token-tới-đầu-ra & $\Theta(T)$ & $\Theta(1)$, kết nối trực tiếp mọi cặp \\ \addlinespace
Chi phí mỗi lớp & $\mathcal{O}(T \cdot d^2)$ & $\mathcal{O}(T^2 \cdot d)$ \\ \addlinespace
Khả năng diễn giải & Khó (trạng thái ẩn ngầm) & Tốt (trọng số chú ý hiển thị) \\ \bottomrule
\end{tabularx}
\end{table}

Từ những phân tích trên, hai hướng nâng cấp kiến trúc tự nhiên được đề xuất nhằm tối ưu hóa hệ thống. Hướng thứ nhất, mang tính tiến hóa, là giữ GRU làm bộ tạo ngữ cảnh nhưng thay thế phép lấy trạng thái ẩn cuối cùng bằng cơ chế chú ý cộng tính (\textit{additive attention pooling}) theo Bahdanau và cộng sự (2015). Cụ thể, với chuỗi trạng thái ẩn $h_1, \dots, h_T$ do GRU sinh ra, ta tính toán trọng số chú ý $\alpha_t$ và tổng hợp thành vectơ ngữ cảnh $c$:
\begin{equation}
e_t = v^\top \tanh(W h_t), \quad \alpha_t = \frac{\exp(e_t)}{\sum_{t'} \exp(e_{t'})}, \quad c = \sum_{t} \alpha_t h_t
\end{equation}
Vectơ ngữ cảnh $c$ sau đó sẽ thay thế hoàn toàn cho $h_T$ để làm đầu vào trực tiếp cho bộ phân loại. Giải pháp này không chỉ dễ tích hợp dựa trên nền tảng sẵn có mà còn tận dụng tốt khả năng nén ngữ cảnh của GRU, đồng thời cung cấp các trọng số $\alpha_t$ trực quan giúp tăng khả năng diễn giải hành vi của mô hình.

Bên cạnh đó, hướng thứ hai mang tính cách mạng hơn là thay thế hoàn toàn mạng hồi quy bằng một cấu trúc \textit{stack Transformer encoder}. Mỗi lớp \textit{encoder} mới này bao gồm một khối \textit{multi-head self-attention} kết hợp cùng mạng \textit{feed-forward} hai tầng, tất cả đều được bọc bởi các kết nối tắt (\textit{residual connection}) và chuẩn hóa tầng (\textit{layer normalization}). Nhằm bổ sung thông tin vị trí cho cơ chế \textit{self-attention}, một bộ mã hóa vị trí (\textit{positional encoding}) hình sin được cộng trực tiếp vào cấu trúc \textit{embedding}:
\begin{align}
\text{PE}(\text{pos}, 2i) &= \sin\left( \frac{\text{pos}}{10000^{2i/d}} \right) \\
\text{PE}(\text{pos}, 2i+1) &= \cos\left( \frac{\text{pos}}{10000^{2i/d}} \right)
\end{align}
Kiến trúc song song hóa toàn diện này được kỳ vọng sẽ giải quyết triệt để các rào cản về mặt tuần tự của GRU, từ đó xử lý hiệu quả các trường hợp châm biếm hoặc phụ thuộc xa như đã ghi nhận trong thực nghiệm.